\section{帰無仮説検定の基本}
\label{b13_NHST}

\subsection{授業内容}
\subsubsection{概要}
前回は統計的推定について学んだ。本コマでは，推定を発展させた統計的意思決定法である帰無仮説検定について学ぶ。帰無仮説検定は心理学研究において広く用いられている統計的推論の方法であり，データに基づいて仮説の妥当性を評価するための枠組みを提供する。本コマでは，帰無仮説検定の基本的な考え方と手順，さまざまな検定法(相関係数の検定，カイ二乗検定など)，そして有意水準や片側・両側検定の概念について理解を深める。帰無仮説検定は背理法を用いて効果の有無を判断する方法であり，結果の報告に関する表記法も含めてその作法を学ぶ必要がある。

\subsubsection{コマ主題細目}
\begin{description}
	\item[統計的帰無仮説検定とその手順] 統計的帰無仮説検定(Null Hypothesis Significance Test; NHST)は帰無仮説に基づいてモデル比較をする一連の手続きである。前回学んだ推定と異なり，検定では特定の仮説(帰無仮説)が正しいと仮定した場合に観測されたデータがどれだけ起こりにくいかを評価する。帰無仮説検定の基本的な考え方，帰無仮説と対立仮説の設定，検定統計量の算出，p値の計算，判断基準の設定という一連の流れを学ぶ。また，検定統計量は標本統計量を標準化したものであり，統計量によって適切な分布(標準正規分布，t分布，カイ二乗分布，F分布など)が対応することを理解する。
	      \begin{flushright}
		      $\to$ 手続きの一般化としては\textcite{Yamada2004}のP.108--109，\textcite{Kawabata201408}のP.81--85.
	      \end{flushright}

	\item[相関係数の検定] 相関係数の統計的検定は，母相関係数がゼロであるという帰無仮説を検証するものである。2変数が無相関(母相関係数$\rho=0$)という帰無仮説の下では，サンプルサイズ$n$と標本相関係数$r$から計算される検定統計量が自由度$n-2$の$t$分布に従うことを理解する。この検定統計量の計算方法と，検定結果の解釈について学ぶ。また，相関係数の検定と同時に，相関係数の信頼区間を求める方法についても触れる。さらに，相関係数の大きさと統計的有意性の違いについても理解を深める。無相関検定はRでは`cor.test()`関数を用いて実行できる。
	      \begin{flushright}
		      $\to$ \textcite{Kawabata201408}のP.98--103，\textcite{kosugi2018}のP.58--60，\textcite{Shimizu2021}のP.150--162.
	      \end{flushright}

	\item[カイ二乗検定] カイ二乗検定は，カテゴリカルデータの分析に用いられる検定法である。観測度数と期待度数の差を評価し，分布の適合度や変数間の独立性を検証する。カイ二乗検定には適合度検定と独立性の検定があり，それぞれの適用場面と計算方法について学ぶ。カイ二乗統計量の算出方法と，カイ二乗分布を用いた確率評価の方法について理解する。また，カイ二乗検定の適用条件(期待度数が5以上など)や結果の解釈についても学ぶ。
	      \begin{flushright}
		      $\to$ カイ二乗検定については\textcite{Kawabata201408}のP.145--153，\textcite{Yamada2004}のP.177--183を参照。
	      \end{flushright}

	\item[有意水準] モデル比較という観点からは，どうしても「判定基準」を考える必要がある。NHSTの文脈では，それは希少性の指標である有意水準になる。有意水準の定め方は領域ごとに定める任意であり，心理学では一般に5\%が用いられている。このことは，複数のモデルの優劣を判定する数字であって，モデルの強さ，正しさを表現する数字ではないことに注意する。有意水準とp値の関係，および有意水準の設定の仕方について理解する。また，p値の正しい解釈についても学び，p値の誤解や誤用についても理解を深める。
	      \begin{flushright}
		      $\to$ \textcite{Yamada2004}のP.112--115，\textcite{Kawabata201408}のP.86--89.
	      \end{flushright}

	\item[片側検定と両側検定] 統計的検定には片側検定と両側検定があり，研究の仮説や関心に応じて適切な方法を選択する必要がある。両側検定は母数が特定の値と異なるかどうかを検証するのに対し，片側検定は母数が特定の値よりも大きい(または小さい)かどうかを検証する。それぞれの場合の帰無仮説と対立仮説の設定方法，臨界値の違い，p値の計算方法，および結果の解釈について理解する。また，検定の方向性を事前に決定することの重要性と，結果を見てから検定方法を変更することの問題点についても学ぶ。
	      \begin{flushright}
		      $\to$ \textcite{Yamada2004}のP.116--119，\textcite{Kawabata201408}のP.90--93.
	      \end{flushright}

\end{description}
\subsubsection{キーワード}
\begin{itemize}
	\item 帰無仮説検定と検定統計量
	\item 相関係数の検定(無相関検定)
	\item カイ二乗検定
	\item 有意水準とp値
	\item 片側検定と両側検定
\end{itemize}
\subsection{授業情報}
\paragraph{コマの展開方法}
適宜投影資料を用いる
\paragraph{標準シラバスにおける位置づけ}
科目番号4；心理学研究法；2.データを用いた実証的な思考方法；C データの統計的記述\\
科目番号5；心理学統計法；1.心理学で用いられる統計手法；F 推測統計:その考え方
\subsubsection{予習・復習課題}
\paragraph{予習}
前回学んだ不偏推定量，点推定と区間推定，標準正規分布の知識を復習しておくこと。特に，標準誤差と信頼区間の概念，および推定と検定の関係について確認しておくと理解が深まる。また，相関係数の概念と計算方法についても再確認しておくこと。
\paragraph{復習}
授業で学んだ内容を定着させるために，以下の点について理解を深めておくこと：

\begin{enumerate}
\item 帰無仮説検定の論理構造と基本的な手順
\item 相関係数の検定の手順と解釈(無相関の検定，検定統計量の計算，結果の解釈)
\item カイ二乗検定の適用場面と使い方(適合度検定と独立性の検定)
\item 有意水準とp値の関係，およびp値の正しい解釈
\item 片側検定と両側検定の違いと，それぞれの適用場面
\end{enumerate}

また，実際に検定を行ってみることで理解が深まる。例えば，相関係数の大きさ，標本の大きさを任意の値にし，無相関検定を実際にやってみると良い。さまざまな数値を入れることで，検定結果がどのように変動するかを理解することは帰無仮説検定の本質を理解するのに役立つ。仮想的な数字では実感が得られにくい場合，身の回りの相関がありそうなデータを具体例として検定し，何が言えるのかを考えてみると良い。